\vspace{-2mm}
\section{Background}
In this section, we provide background on the properties and existing approaches to support agentic inference.
% Additional discussion of related work is in \textcolor{red}{X}.
% \vspace{-2mm}
\subsection{System Properties of Current Agentic Workflows}
\label{sec::react agent system}

Current agentic workflows alternate between reasoning and acting during generation. Formally, at each step $t$, the agent receives an observation $o_t \in \mathcal{O}$ and produces an emission $e_t = (\ell_t, a_t) \in \mathcal{L} \times \mathcal{A}$, where $\ell_t$ denotes a thought and $a_t$ represents an action. We define the cumulative context at step $t$ as $c_t = (o_1, e_1, \dots, o_t)$, which captures the interaction history of agentic workflows. Conditioned on $c_t$, $e_t$ is sampled from a policy $\pi(e_t|c_t)$.

This workflow keeps two persistent states: \textit{(i) GPU Memory}, where the KV cache of $c_t$ serves as the workflow’s memory. As the trace grows incrementally, $c_{t+1}$ extends $c_t$ as a prefix, enabling theoretical near-complete KV cache reuse rates across steps. \textit{(ii) Tool Environment}, where external resources (e.g., sandboxes or database connections) initialized at $t=1$ \mbox{must remain consistent and accessible throughout the execution.}

These stateful dependencies necessitate a \textit{program-level} view of agentic inference trajectories, thereby enabling to system to coordinate heterogeneous resources and manage state across long-running workflows. However, existing inference systems treat each thought $l_t$ and action $a_t$ as an independent, stateless request.

\vspace{-2mm}
\subsection{Existing Agentic Inference Systems}

Prior work focuses on optimizing the individual components in agentic inference, including the LLM inference engine or tool orchestrator ~(\autoref{appendx: kv opt}, \autoref{sec: pd+offload}, but there are very few works that provide end-to-end optimization for agentic workflows across GPU, CPU, and remote resources. We review these prior systems.
% fail to maintain high throughput in agentic serving and RL rollout discussed in \autoref{appendx: kv opt} and \autoref{sec: pd+offload}
% as a heterogeneous component.

Autellix models multi-turn agentic workflows as \textbf{GPU-only programs} and tracks the accumulated GPU execution time in a central process table~\citep{luo2025autellixefficientservingengine}. However, it ignores workflow locality, allowing concurrent workflows to aggressively evict other's KV cache, triggering \textbf{KV cache thrashing} under heavy workloads.


Continuum is another recent serving system designed for multi-turn agentic workflows~\cite{li2025continuumefficientrobustmultiturn}. It employs a time-to-live (TTL) mechanism to pin KV caches in HBM, thereby mitigating context thrashing during tool execution.
However, it fails to solve the KV cache eviction problem. The first reason is that most tools take an \textbf{unpredictable} amount of time (e.g. remote model APIs in ToolOrchestra~\citep{su2025toolorchestraelevatingintelligenceefficient}, compilers in code agents, and web applications for computer use agents~\citep{zhou2024webarenarealisticwebenvironment}).
Such unpredictable tools trigger severe thrashing as well as stranded KV cache memory in Continuum due to incorrect TTL estimates.
Moreover, once the decoding memory of the running workflow surpasses the GPU limit, the system preempts and evicts pinned KV cache as well. This leads to unavoidable thrashing and corresponding throughput degradation shown in \autoref{fig:throughputs_vs_bs}.

% These limitations necessitate a unified and protable system for agentic workloads, serving as a plug-and-play layer for agentic infrastructures.

These limitations underscore the need for a simple and fast system for agentic inference. We envision such a system as a program-aware scheduling layer for emerging agentic inference systems (e.g., \citet{zhang2026megaflow}).
% , unlocking scalable serving and RL efficiency.